% =====================================================================
%  Resume et abstract, en quatrieme de couverture
% =====================================================================
\cleardoublepage
\thispagestyle{empty}
\vspace*{1cm}

\begin{center}
  {\Large\sffamily\bfseries\color{bleuTitre} Resume\par}
\end{center}
\vspace{0.5cm}

Le transport routier de marchandises repose sur la coordination de
quatre ressources aux cycles de vie independants : un parc de vehicules,
des equipements embarques, des chauffeurs et un portefeuille de clients.
Les entreprises de taille moyenne gerent cette coordination avec des
outils qui n'y sont pas destines, ce qui produit une information
dispersee, une maintenance subie et une relation client opaque.

Ce projet de fin d'etudes, realise au sein de \entreprise, a consiste a
concevoir et developper \sujetCourt, une plateforme de gestion de flotte
et de transport de marchandises fondee sur une architecture
microservices. Elle reunit huit services metier en .NET~8 appliquant la
Clean Architecture, quatre services d'aide a la decision en Python, et
trois applications destinees a des publics distincts : un back-office, un
portail client et une application mobile pour les chauffeurs. Le tout est
deploye en vingt conteneurs orchestres par Docker Compose.

Le travail a ete conduit selon le cadre Scrum, en six sprints couvrant
successivement les identites et les acces, le referentiel de flotte, la
chaine commerciale de la demande au suivi en direct, la maintenance
curative et predictive, la facturation et les alertes, enfin la
collaboration et l'aide a la decision.

\vspace{0.4cm}
\noindent
\textbf{Mots-cles :} microservices, gestion de flotte, transport de
marchandises, Clean Architecture, Scrum, aide a la decision.

\vspace{1.4cm}
{\color{grisFilet}\rule{\textwidth}{0.8pt}}
\vspace{1.4cm}

\begin{center}
  {\Large\sffamily\bfseries\color{bleuTitre} Abstract\par}
\end{center}
\vspace{0.5cm}

Road freight transport relies on coordinating four resources with
independent life cycles: a vehicle fleet, on-board equipment, drivers and
a customer portfolio. Medium-sized companies handle this coordination
with tools that were never designed for it, which leads to scattered
information, reactive maintenance and limited visibility for customers.

This final-year project, carried out at \entreprise, consisted in
designing and developing \sujetCourt, a fleet and freight transport
management platform built on a microservices architecture. It brings
together eight .NET~8 business services following Clean Architecture,
four Python decision-support services, and three applications serving
distinct audiences: a back-office, a customer portal and a mobile
application for drivers. The whole system is deployed as twenty
containers orchestrated by Docker Compose.

The work followed the Scrum framework, over six sprints covering
identity and access management, the fleet repository, the commercial
chain from request to live tracking, corrective and predictive
maintenance, invoicing and business alerts, and finally collaboration and
decision support.

\vspace{0.4cm}
\noindent
\textbf{Keywords:} microservices, fleet management, freight transport,
Clean Architecture, Scrum, decision support.
